\SourceLeaf{378}
\par\bigskip\Needspace{9\baselineskip}
\pdfbookmark[1]{Theory of Biquadratic Residues. VI.}{nachlass-biquadratic-VI}
\centerline{[VI.]}\medskip
\begin{center}\large\scshape Theory of Biquadratic Residues.\end{center}
\BookRule{.18\textwidth}
\renewcommand{\Article}[1]{\par\medskip\Needspace{4\baselineskip}\pdfbookmark[2]{Art. #1}{nachlass-biquadratic-VI-art-#1}\centerline{\textbf{#1.}}\nobreak\smallskip}
\Article{1}
Let a sequence of complex integers $\varphi,\varphi',\varphi''$, and so on, satisfy two conditions: first, they are all pairwise incongruent modulo $\mu=\alpha+\beta i$, where $\alpha$ and $\beta$ denote real integers; second, every complex integer is congruent to one of them modulo $\mu$. Under this assumption, the collection of numbers $\varphi,\varphi',\varphi''$, and so on, is called a complete residue system for the modulus $\mu$. It has been proved that the number of its members equals the norm of $\mu$, that is, the number $\alpha\alpha+\beta\beta$, which will be denoted by $\nu$.

\Article{2}
Among the numbers $\varphi,\varphi',\varphi''$, and so on, one is divisible by $\mu$. If it is excluded and the collection of the others is denoted by $\chi$, then $\chi$ forms a complete system of residues not divisible by the modulus, their number being $=\nu-1$. If we restrict the investigation to odd moduli, $\nu-1$ is divisible by $4$. Furthermore, $f,if,-f,-if$ will then be pairwise incongruent; hence the numbers in $\chi$ congruent, respectively, to $if,-f,-if$ are distinct from one another and from $f$. (Associated and composite numbers.)

This partitions $\chi$ into four groups or partial systems $\kappa,\kappa',\kappa'',\kappa'''$. Place any number from $\chi$, for example $\varphi$, in the group $\kappa$, and the three members of $\chi$ congruent to $i\varphi,-\varphi,-i\varphi$ in the groups $\kappa',\kappa'',\kappa'''$, in that order. After removing these four numbers from $\chi$, again place any one of the remaining numbers in $\kappa$, and the three that depend on it in the same way in $\kappa',\kappa'',\kappa'''$. Continue until the entire system $\chi$ has been distributed. The groups $\kappa,\kappa',\kappa'',\kappa'''$ will be called associated quarter systems. They clearly have the following properties:

\SourceLeaf{379}
\begin{enumerate}
\renewcommand{\labelenumi}{\arabic{enumi})}
\item Each quarter system consists of $\frac14(\nu-1)=\frac14(\alpha\alpha+\beta\beta-1)$ numbers.
\item The characteristic property of a quarter system is that none of its numbers, either itself or multiplied by $i,-1$ or $-i$, is congruent to another number in the same quarter system, whereas every number not divisible by $\mu$, either itself or multiplied by $i,-1$ or $-i$, occurs in it or is congruent to one of its members.
\item Just as multiplying the numbers in $\kappa$ by $i,-1$ and $-i$ gives, respectively, the numbers in $\kappa',\kappa'',\kappa'''$, or numbers congruent to them, multiplying the numbers in $\kappa'$ by those factors reproduces, respectively, the numbers in $\kappa'',\kappa''',\kappa$; multiplying the numbers in $\kappa''$ similarly reproduces the numbers in $\kappa''',\kappa,\kappa'$; finally, multiplying the numbers in $\kappa'''$ reproduces $\kappa,\kappa',\kappa''$.
\end{enumerate}

For brevity, this mutual dependence of the four quarter systems can be expressed as
\[
\begin{aligned}
\kappa'&\equiv i\kappa,\\
\kappa''&\equiv-\kappa\equiv i\kappa',\\
\kappa'''&\equiv i\kappa''\equiv-\kappa'\equiv-i\kappa.
\end{aligned}
\]
